\begin{frame}
    \frametitle{Conclusion}
    \begin{itemize}
        \item $k$-reducibility extends upon ideas from the Five Color Theorem.
        \item D-reducibility extends upon $k$-reducibility for rings 6 and above.
        \item C-reducibility extends upon D-reducibility by avoiding bad ring colorings.
        \item Unavoidability guarantees that every planar graph has a reducible configuration.
        \item The Birkhoff Diamond and Bernhart Diamond illustrate D and C-reducibility.
        \item Symmetry faults in the Bernhart Diamond require more attention.
    \end{itemize}
\end{frame}

\begin{frame}
    \frametitle{Conclusion}

    The more advanced the concept of reducibility, the less reducible configurations are needed (A/B-reducibility is a case of C reducibility).

    \vspace{1cm}
    \begin{figure}
        \centering
        \includegraphics[width=1.0\textwidth]{images/reduce.eps}
        \caption{Number of reducible configurations on rings $R_6$ and higher in proofs of the four color theorem}.
    \end{figure}

    \uncover<2->{
        Does "Perfect reducibility" exist? If so, it will be the heart of the four color theorem.
    }
\end{frame}

\begin{frame}
    \frametitle{The End}
    Thank you for your attention. Questions?
\end{frame}